\section{Threat model and problem statement}
\label{sec:problem}

\subsection{Setting}

A PDP receives a stream of access requests $e_1, e_2, \dots$ in arrival order. Request
$e_i$ at time $t_i$ carries a client network origin $s_i$, a TLS configuration
fingerprint $c_i$, a device identity $d_i$, an authenticated principal $u_i$, a
requested resource $r_i$, and a feature vector $m_i \in \mathbb{R}^{10}$ (request
method, sensor state, TLS-trust bit, request volume, inter-arrival gap). The detector
must emit an anomaly score $a_i \in [0,1]$ using strictly $e_1, \dots, e_i$, without access
to future events or batch-level lookahead, matching operational streaming conditions.

Two direct consequences follow. First, response-side fields are excluded from
$m_i$, because HTTP response status codes and byte sizes are unavailable when authorization
decisions occur; empirical validation confirmed that including response features introduces
deterministic label leakage. Second, the evaluation replays the test stream event by
event, strictly scoring each event before its memory state is updated.

\subsection{Attacker}

The adversary obtains valid credentials (e.g., via credential harvesting, host compromise,
or session theft) and issues requests that the policy engine permits. The attacker may
operate custom tooling from an external machine and network, and may perform reconnaissance
prior to moving laterally. The adversary is assumed not to have simultaneously replicated
the victim's TLS configuration, device attestation, and behavioral history, as discussed
in Section~\ref{sec:limits}.

Access requests are categorized into four classes:

\begin{itemize}
  \item \textit{Policy violation}: Role, clearance, or sensitivity constraints are unsatisfied.
  These requests are rejected upstream by the policy engine.
  \item \textit{Contextual}: Requests exhibiting invalid TLS posture or deterministic
  intrusion alerts, which are predominantly handled by static rule engines.
  \item \textit{Lateral movement}: Authorized but non-habitual requests that are feature-identical
  to legitimate traffic, forming the primary objective of temporal graph modeling.
  \item \textit{Credential reuse}: Valid credentials presented from a client TLS configuration
  never previously associated with that principal, identifiable via the
  \textit{config}\,$\rightarrow$\,\textit{user} binding.
\end{itemize}

The 5-node graph architecture is designed to capture these latter two classes, for which
the leakage audit described in Section~\ref{sec:setup} strictly applies.

\subsection{Learning regime}

Training follows a \textit{one-class} semi-supervised regime. Labels are used exclusively
in two stages: to filter the training set (benign access events only) and to calibrate the
decision threshold on a validation slice. No ground-truth label enters test-time inference.
The primary operational threshold (persisted and served) is the \textit{cost-sensitive}
threshold, calibrated with labeled lateral movement events in the validation window. A
label-free variant (a quantile of the benign score distribution at a target false-positive
rate) is recorded in the calibration metadata for deployments where red-team labels are
unavailable.

\subsection{The commit gate}

\begin{figure}[ht!]
\centering
\begin{tikzpicture}[
  font=\sffamily\scriptsize,
  box/.style={draw, rounded corners=2pt, align=center, inner sep=3.5pt, minimum height=6mm},
  fl/.style={->, >=Stealth, semithick}
]
\node[box, fill=blue!8] (ev) {access request $e_i$};
\node[box, fill=orange!10, below=5mm of ev] (score) {score $a_i$ \\ \textit{(memory read only)}};
\node[box, fill=green!10, below=5mm of score] (opa) {policy engine};
\node[box, fill=gray!10, below left=6mm and -6mm of opa] (upd) {commit: \\ memory $+$ neighbours};
\node[box, fill=red!8, below right=6mm and -6mm of opa] (noupd) {\textbf{no} commit};
\draw[fl] (ev) -- (score);
\draw[fl] (score) -- node[right, font=\sffamily\tiny] {$a_i$ advisory} (opa);
\draw[fl] (opa) -- node[left, font=\sffamily\tiny, pos=0.6] {ALLOW} (upd);
\draw[fl] (opa) -- node[right, font=\sffamily\tiny, pos=0.6] {DENY} (noupd);
\end{tikzpicture}
\caption{The streaming decision flow. The model scores against memory but does not govern
what enters it; the external policy engine does. The score is returned as an advisory signal
and can be overridden by external policy rules.}
\label{fig:gate}
\end{figure}

A streaming detector maintaining a dynamic baseline faces potential self-poisoning: if the
model ingests traffic based solely on its own classification, an adversary can gradually
shift the baseline. Decoupling inference from memory updates (Fig.~\ref{fig:gate})
addresses this: entity memory and temporal neighborhoods advance only on requests
admitted by the external policy engine, preventing error compounding and runaway
baseline drift. All evaluations in this paper operate under this external commit protocol.

In the streaming evaluation, this commit gate evaluates a deterministic function
over edge signals inspected by an external policy engine (such as invalid TLS posture or
sensor alerts), implemented as an explicit admission filter (\texttt{signal\_dirty}) in
\texttt{serve\_tgn.py} and \texttt{train\_tgn.py}. Evaluating against this deterministic signal
establishes a conservative baseline: authorized lateral movement events with clean signals
enter memory, exposing the detector to realistic stream poisoning during evaluation.
Direct integration with an interactive external policy service (\texttt{policy\_allows},
evaluated in \texttt{tests/test\_policy\_blp.py}) serves as the reference architecture.

Two operational considerations follow from this design: authorized adversary activity that
passes policy checks will still update entity memory, and benign entities persistently flagged
as anomalous risk starvation if denied updates. In practice, authorization frameworks address
these trade-offs through strict organizational boundary policies and warm-up grace periods
for newly provisioned entities.
